\section{自动微分给出导数，梯度检查建立信任}
\label{sec:08-Numerical-Analysis/07-Numerical-Methods-for-ML/02}

你给模型加了一个自定义层，前向照着公式写，反向自己推了一遍。训练跑起来了，损失在降——降得比预期慢一点，但确实在降。

这是最难查的一类故障。梯度里少了一项、转过了一次、广播错了一个轴，模型往往还是能学到东西，因为剩下的那部分梯度方向大致对头。没有任何断言会失败，只有一条比参考实现平坦一点的曲线。

要判断这段求导代码写没写错，需要一条独立的计算路径给出同一个答案。最容易想到的是差商：把某个参数挪一点，看损失变多少。

\[
f'(x)\approx\frac{f(x+h)-f(x)}{h}.
\]

这条路径确实独立，代价却大得没法用来训练：一百万个参数就要重跑一百万遍模型。而且它自己也不准——\hyperref[sec:08-Numerical-Analysis/04-Numerical-Differentiation-and-Integration/01]{``数值微分为何不能无限缩小步长''}把这条权衡算得很细：截断误差随 \(h\) 减小而降，舍入误差随 \(h\) 减小而升，总误差是一条 U 形曲线，谷底处中心差分也只剩十位左右的有效数字。

两件事因此要分开：训练要的是高效而精确的导数，验证要的是一条与被验证代码尽可能无关的路径。前者归自动微分，后者归差商。它们不是竞争关系，把它们摆在一起比较准不准，是很多困惑的来源。

\subsection{自动微分对程序求导，不对公式求导}

自动微分既不是差商，也不是把表达式展开成符号形式再化简。它把程序拆成一串基本运算，每个运算的局部导数是已知的，然后沿着这些运算构成的图传播。

拿

\[
y=\sin(x_1x_2)+x_2^2
\]

来说，程序实际执行的是 \(a=x_1x_2\)，\(b=\sin a\)，\(c=x_2^2\)，\(y=b+c\)。这四步和它们之间的依赖关系就是计算图。

\begin{center}
\begin{tikzpicture}[x=1cm,y=1cm,font=\small,>=stealth]
  \draw[fill=black!6,draw=black!50] (0,2.6) circle (0.42);
  \node[font=\footnotesize] at (0,2.6) {\(x_1\)};
  \draw[fill=black!6,draw=black!50] (0,1.0) circle (0.42);
  \node[font=\footnotesize] at (0,1.0) {\(x_2\)};
  \draw[fill=blue!12,draw=black!50] (2.4,2.6) circle (0.42);
  \node[font=\footnotesize] at (2.4,2.6) {\(a\)};
  \draw[fill=blue!12,draw=black!50] (4.8,2.6) circle (0.42);
  \node[font=\footnotesize] at (4.8,2.6) {\(b\)};
  \draw[fill=blue!12,draw=black!50] (2.4,0.6) circle (0.42);
  \node[font=\footnotesize] at (2.4,0.6) {\(c\)};
  \draw[fill=orange!18,draw=black!50] (7.2,1.6) circle (0.42);
  \node[font=\footnotesize] at (7.2,1.6) {\(y\)};
  \draw[->,black!60] (0.42,2.6) -- (1.98,2.6);
  \draw[->,black!60] (0.35,1.23) -- (2.05,2.37);
  \draw[->,black!60] (0.41,0.93) -- (1.99,0.67);
  \draw[->,black!60] (2.82,2.6) -- (4.38,2.6);
  \draw[->,black!60] (5.19,2.44) -- (6.81,1.76);
  \draw[->,black!60] (2.81,0.69) -- (6.79,1.51);
  \node[above,font=\footnotesize,black!65] at (1.2,2.62) {\(x_2\)};
  \node[font=\footnotesize,black!65] at (1.55,1.9) {\(x_1\)};
  \node[below,font=\footnotesize,black!65] at (1.2,0.72) {\(2x_2\)};
  \node[above,font=\footnotesize,black!65] at (3.6,2.62) {\(\cos a\)};
  \node[above,font=\footnotesize,black!65] at (6.1,2.2) {\(1\)};
  \node[below,font=\footnotesize,black!65] at (4.9,1.02) {\(1\)};
  \draw[->,blue!70] (-0.4,3.6) -- (7.6,3.6);
  \node[above,font=\footnotesize,blue!75] at (3.6,3.62) {正向模式：与执行同向，一次给出 \(Jv\)};
  \draw[->,red!70] (7.6,-0.5) -- (-0.4,-0.5);
  \node[below,font=\footnotesize,red!75] at (3.6,-0.52) {反向模式：逆着边回溯，一次给出 \(v^{\trans}J\)};
\end{tikzpicture}

\emph{每条边上标的是一步基本运算的局部导数。两种模式走的是同一张图、用的是同一批局部导数，区别只在遍历方向：一个从输入往输出推，一个从输出往输入收。}
\end{center}

正向模式给每个变量额外挂一个切量，与数值同步往前推：给定输入方向 \(v\)，一趟下来得到 Jacobian 与该方向的乘积 \(Jv\)，也就是所有输出对这一个输入方向的敏感度。反向模式在输出端放一个伴随量再逆着边走，一趟给出 \(v^{\trans}J\)，也就是一个输出对所有输入的梯度。图里那两条粗箭头标的就是这件事。

反向的账在例子上过一遍。记 \(\bar z=\partial y/\partial z\)，从 \(\bar y=1\) 出发，加法把它原样分给两支，\(\bar b=\bar c=1\)；经过 \(b=\sin a\) 得 \(\bar a=\cos a\)；\(\bar a\) 沿 \(a=x_1x_2\) 按局部导数 \(x_2\) 与 \(x_1\) 分给两个输入，\(\bar c\) 沿 \(c=x_2^2\) 以 \(2x_2\) 流向 \(x_2\)，在 \(x_2\) 处两股汇合相加：

\[
\begin{aligned}
\frac{\partial y}{\partial x_1}&=x_2\cos(x_1x_2),\\
\frac{\partial y}{\partial x_2}&=x_1\cos(x_1x_2)+2x_2.
\end{aligned}
\]

得到的正是手推公式，没有任何近似——图上每条边的局部导数是解析的，链式法则是精确的，唯一的误差来自那些乘法和加法本身的舍入。差商那条 U 形曲线所权衡的截断误差，在这里根本不存在。这就是自动微分能同时做到``快''和``准''的地方。

\subsection{方向由输入输出的维数定}

对 \(f:\R^n\to\R^m\)，完整 Jacobian 有 \(mn\) 个元素，两种模式都不直接构造它。正向模式一趟给一列方向的 \(Jv\)，凑齐整个 Jacobian 要 \(n\) 趟；反向模式一趟给一行的 \(v^{\trans}J\)，凑齐要 \(m\) 趟。哪个便宜，看 \(n\) 和 \(m\) 谁小。

神经网络训练把这个取舍推到了极端：\(m=1\)（一个标量损失），\(n\) 是上亿个参数。反向模式一趟就够，代价约为一次前向的两三倍，与参数个数无关。这就是反向传播——它不是深度学习发明的一套特殊算法，而是反向模式自动微分用在标量损失上的名字。相反，若要考察一个参数的扰动如何影响整条轨迹的所有状态（\(n=1\)、\(m\) 很大），正向模式才是省的那个。

反向模式省时间，省下的时间要用空间来付：局部导数要用到前向的中间值，所以整张图的中间量得留到反传时用。显存不够时的常规做法是重计算，前向只存少数检查点，反传路过时再把中间段重新算一遍，拿时间换空间。这条取舍在\hyperref[sec:14-Deep-Learning/02-Losses-and-Backpropagation/02]{``一次完整的前向与反向传播''}中有具体的层级展开，图与伴随量的一般写法见\hyperref[sec:10-Matrix-Calculus/04-Applications/03]{``计算图上的反向传播与自动微分''}。

还有一件常被忽略的事：许多任务要的从来不是 Jacobian，而是它对某个向量的作用。敏感度分析要 \(Jv\)，标量损失反传要 \(J^{\trans}v\)，Newton 型方法要的 Hessian-向量积可以由两种模式套用得到，隐式层的反传只需解一个伴随线性系统。先分清要的是整个矩阵还是矩阵乘向量，能省掉大量时间和存储——这与下一节\hyperref[sec:08-Numerical-Analysis/07-Numerical-Methods-for-ML/03]{``需要的是解，而不是逆矩阵''}是同一条判断的两次应用。

\subsection{精确求导，不等于求出稳定的导数}

这里的``精确''需要限定：自动微分严格按照\emph{所执行的那个程序}套用链式法则，它不会替你换一条更好的计算路径。前向计算里若有灾难性抵消、有溢出、有一个已经变成 NaN 的中间量，导数会忠实地继承这一切，有时还放大。上一节那个先 softmax 后取对数的写法，用自动微分求导，得到的梯度同样是 NaN。

顺序因此是固定的：先把前向改写稳定，再谈求导。反过来指望求导环节把一个不稳定的前向救回来，不会发生。

\subsection{一次检查要能失败，才算建立了信任}

设自动微分给出 \(g_{\mathrm{AD}}\)。逐坐标比较用中心差分，

\[
g_i^{\mathrm{num}}=\frac{f(x+he_i)-f(x-he_i)}{2h},
\]

参数一多就跑不动。更经济的是投影检查：随机取一个方向 \(v\)，比较 \(\big(f(x+hv)-f(x-hv)\big)/(2h)\) 与 \(v^{\trans}g_{\mathrm{AD}}\)。一次投影把所有分量混在一起，能迅速抓到整体性的错误；它的盲区是那些几乎与 \(v\) 正交的误差，所以要换几个方向，并在能承受的小规模问题上补一次逐坐标检查。

几个细节决定这次检查有没有判断力。用中心差分而不是前向差分，截断误差从 \(O(h)\) 降到 \(O(h^2)\)，谷底更低更平；\(h\) 取中间量级，float64 下 \(10^{-5}\) 到 \(10^{-6}\) 附近通常合适，并且要按变量自身的尺度缩放，不能对量级为 \(10^{6}\) 的参数和量级为 \(10^{-3}\) 的参数用同一个绝对步长；比较相对误差而非绝对差，常用的尺度化指标是

\[
\frac{\norm{g_{\mathrm{AD}}-g_{\mathrm{num}}}}{\max\set{1,\norm{g_{\mathrm{AD}}},\norm{g_{\mathrm{num}}}}},
\]

分母里那个 \(1\) 是为了梯度本身接近零时不把噪声放大成失败。阈值没有脱离上下文的统一答案，float64 上通过的量级到了 float32 上一定不成立。最后，别只试一个 \(h\)：扫过几个数量级，看误差是否先降后升、谷底是否落在预期位置。谷底比预期高一大截，是实现有错的强信号；曲线根本没有谷底，多半是测试点本身有问题。

测试环境也要配合：输入小而可读，尽量用双精度；固定随机种子，关掉 dropout，冻结批归一化那类会随调用更新的统计量；避开 ReLU 的零点、\(\max\) 的并列点和 clip 的边界。

最后一步最容易被跳过，却最重要——故意注入一个错误，看检查会不会红。把某个梯度乘上 \(-1\)、去掉一次转置、把一个求和的轴换掉，如果检查照样通过，那它此前的绿灯什么也没证明。

\subsection{不光滑处的分歧未必是错}

\(f(x)=\abs{x}\) 在 \(x=0\) 处没有导数。对称差分给出 \(0\)，多数自动微分系统按约定也返回 \(0\)，两者一致，但这种一致不代表它们找到了一个经典意义上的导数——它只是两套约定碰巧撞在一起。带数据依赖分支、离散采样、排序或原地状态更新的程序，情况更微妙，常常需要自定义求导规则，或者把模型重新表述成处处可微的形式。

梯度检查回答的问题因此有明确边界：在选定的光滑测试点上、在确定的程序状态下、按既定的求导约定，两条独立路径给出的数是否一致。它不回答模型是否可微，那是写模型时就该判断清楚的事。

导数拿到了，接下来多半要拿它去解一个方程——下一节说的就是这一步里最常见的那个误译。
